\section{与 EP139 的对照：技术史框架到实验室行动}

EP139 苏煜访谈给出 Agent 技术史：Agent 作为实体、环境和目标导向行动的系统，如何从 logical agent、semantic parsing 走到 language agent 和 universal digital agent。EP138 罗福莉访谈则回答另一个问题：当这个范式来到模型实验室内部，团队应该如何行动。

\noindent\begin{tabularx}{\textwidth}{p{0.22\textwidth}p{0.34\textwidth}X}
\toprule
维度 & EP139 苏煜 & EP138 罗福莉 \\
\midrule
关注层级 & Agent 技术谱系与概念框架 & AI Lab 的资源、组织和模型路线。 \\
核心转折 & OpenClaw Moment 让 Agent 共识形成 & OpenClaw 改变团队研究范式。 \\
关键机制 & Language Agent、边界消融、continual learning & Post-train、RL infra、多模型编排、卡分配。 \\
落地问题 & reliability、cost、deployment、world model & 端到端完成率、成本速度、组织平权、环境反馈。 \\
\bottomrule
\end{tabularx}

\begin{knowledgebox}{连读方法}
先读 EP139，理解 Agent 为什么是一个长期技术问题；再读 EP138，理解模型团队如何在这个技术问题变成产业共识后调整资源、训练系统和组织结构。
\end{knowledgebox}

\subsection{本章小结}

两期合起来构成一组：EP139 是 Agent 的概念地图，EP138 是 Agent 时代模型实验室的行动地图。后续Zhang Xiaojun AI 队列可以沿着这两张地图继续整理。


